\chapter{Application Threat Auditing}
\label{ch:5.6}
\chaptertopic{Topic 5.6}

\chapterquote{Trust, but verify.}{Russian proverb, popularized in English by Ronald Reagan}

\section{Motivation}

Chapter~\ref{ch:5.5} built defenses into the application itself. This final chapter of the course
is about what happens once an attack slips past those defenses anyway. It covers how an
application's own logs reveal it, how data leaving the organization is caught in the act, and how
scattered evidence becomes a defensible conclusion about what actually happened. It closes the same
loop Chapter~\ref{ch:4.4} opened for devices, applied here to applications and data. This is the
final layer of the course, and as always we pair an attack with the
specific control that answers it.

\begin{apalignment}{Unit 5: Securing Applications and Data}{Topic 5.6 -- Detecting Attacks on Data and Applications}
\begin{apobjectives}
\apobjective{5.6.A}{Detect Attacks}{Account for how an attack against stored data becomes visible at all -- which record of the event exists, which system wrote it, and what an analyst must compare it against before it means anything.}
\apobjective{5.6.B}{Detect Attacks}{Choose the controls that will actually produce that evidence for a given application and a given threat, rather than assembling a collection of tools that all watch the same thing.}
\apobjective{5.6.C}{Detect Attacks}{Judge a detection method on what it can and cannot see, the noise it generates per true finding, how quickly it fires, what it costs to operate, and whether an adversary can disable it.}
\apobjective{5.6.D}{Detect Attacks}{Determine whether a file is the file it should be by computing its digest and comparing that against a trusted reference, and recognize when the reference is not trustworthy.}
\apobjective{5.6.E}{Detect Attacks}{Read a sequence of log records well enough to name the attack, assess whether it succeeded, and write the result as a report someone else can act on without repeating the analysis.}
\end{apobjectives}
\end{apalignment}

\section{Application Monitoring and Logging}

Every login, form submission, and error an application handles can be recorded as a structured,
timestamped log line. An attacker probing for valid usernames, hammering a login form, or calling an
endpoint they should not reach touches the application in exactly the same way a legitimate user
does. That means every one of those actions can leave a mark.

\begin{example}[The pattern is the signal, not the event]
{\scriptsize
\begin{verbatim}
08:42:01Z  192.168.10.44  jsmith  POST /login  401  "Invalid credentials"
08:42:03Z  192.168.10.44  jsmith  POST /login  401  "Invalid credentials"
08:42:05Z  192.168.10.44  jsmith  POST /login  200  "Login successful"
\end{verbatim}
}
Seen once, this is unremarkable -- people mistype passwords. Seen fifty times, from a single IP,
across twenty usernames within two minutes, it is a brute-force attack identical in shape to the
one detected in Chapter~\ref{ch:4.4}. A single log line rarely proves anything; the correlated
pattern across many lines is what actually carries evidentiary weight.
\end{example}

Several recurring signatures recur across real incidents. \textbf{Brute-force / credential
stuffing} shows up as a high volume of 401 or 403 responses to a login endpoint from one IP,
targeting many usernames in a short window. \textbf{SQL injection probes} appear as URL parameters
or body fields containing the SQL meta-characters from Chapter~\ref{ch:5.1} -- \texttt{'},
\texttt{--}, \texttt{UNION SELECT} -- captured verbatim in a raw access log. \textbf{Privilege
escalation attempts} show up as a series of 403 responses from an authenticated low-privilege
account suddenly requesting \texttt{/admin/} or a configuration endpoint. \textbf{Error storms} --
a sudden spike in 500-series responses concentrated on one endpoint -- can indicate active
exploitation of a code vulnerability forcing the server into an unexpected state.

\begin{lstlisting}[language=Python, caption={Flagging IPs exceeding a brute-force threshold from parsed log entries}]
from collections import defaultdict

def flag_bruteforce_ips(entries, threshold=5, window_seconds=60):
    """entries: list of dicts like {"ip": str, "timestamp": int,
       "status": int}. Returns sorted IPs with >= threshold 401/403
       responses within any window_seconds window."""
    by_ip = defaultdict(list)
    for e in entries:
        if e["status"] in (401, 403):
            by_ip[e["ip"]].append(e["timestamp"])

    flagged = set()
    for ip, times in by_ip.items():
        times.sort()
        left = 0
        for right in range(len(times)):
            while times[right] - times[left] > window_seconds:
                left += 1
            if right - left + 1 >= threshold:
                flagged.add(ip)
                break
    return sorted(flagged)
\end{lstlisting}

Security-relevant events must be logged at INFO severity or higher, since a misconfiguration that
drops them creates a blind spot an attacker can exploit by deliberately staying quiet enough to
avoid it. Because raw logs are high-volume -- a busy server can produce thousands of lines per
second -- teams feed them into a SIEM, exactly as in Chapter~\ref{ch:3.5}. SIEM rules apply logic
like ``if one IP produces five or more 401s within sixty seconds, alert.'' Logging is not complete
protection on its own: it cannot capture traffic that bypasses the application entirely, such as a
compromised insider reading the database directly. An attacker who achieves code execution will
often try to erase logs to cover their tracks. That is exactly why logs should stream to a separate,
append-only, remote aggregator in real time, so local deletion does not erase the forensic record.

\begin{misconception}
\textbf{``If no alarm fired, nothing bad happened.''} Alarms only fire on configured rules; a
missing or misconfigured rule is a silent blind spot regardless of what actually occurred. Periodic
manual log review remains necessary even with automated alerting in place.
\end{misconception}

\section{Data Loss Prevention}

An organization can defend its perimeter, patch its applications, and monitor its logs, and still
lose its most valuable asset quietly: the data itself.

\begin{definition}{Data loss prevention (DLP)}
Tools and processes that detect, and where possible block, sensitive data leaving an organization's
control without authorization -- whether the cause is a malicious insider, a compromised account,
or an honest employee's mistake.
\end{definition}

Two detection strategies work together. \textbf{Content-based detection} inspects the actual data
leaving the organization for recognizable patterns -- a credit-card number format, an SSN pattern, a
document labeled ``CONFIDENTIAL'' in its metadata. It is precise when it fires, but blind to
anything an attacker compresses or encrypts before sending, since the recognizable pattern simply
disappears. \textbf{Volume/behavior-based detection} looks instead at the \emph{shape} of data
movement rather than its content: an account that normally transfers a few megabytes a day suddenly
uploading two gigabytes to an unfamiliar destination is anomalous regardless of what that data
actually contains. That is exactly what lets behavior-based detection catch encrypted exfiltration
that content inspection cannot see at all.

\begin{lstlisting}[language=Python, caption={Flagging outbound transfers by volume and destination}]
def flag_exfiltration(transfers, approved_destinations, volume_threshold_mb=500):
    """transfers: list of dicts like {"account": str, "destination": str,
       "volume_mb": int}. Returns sorted (account, destination) pairs
       exceeding volume_threshold_mb to a non-approved destination."""
    flagged = []
    for t in transfers:
        if (t["destination"] not in approved_destinations and
                t["volume_mb"] > volume_threshold_mb):
            flagged.append((t["account"], t["destination"]))
    return sorted(flagged)
\end{lstlisting}

A mature DLP program instruments multiple vantage points, since different exfiltration paths are
visible from different places. \textbf{Network DLP} watches traffic at the perimeter.
\textbf{Endpoint DLP} runs on individual devices and can catch a USB copy that never touches the
network at all. \textbf{Cloud/email DLP} inspects uploads and attachments before they leave a
managed environment. Because legitimate large transfers happen constantly -- backup jobs, sanctioned
partner data-sharing -- effective DLP maintains an \textbf{approved destination list}, so an
identical transfer volume is silent to a known-good partner and alarming to an unrecognized external
domain. Building and tuning that list is most of the ongoing work of actually running a DLP program.

\begin{example}[Why DLP is preventive, not merely forensic]
Unlike a log reviewed after the fact, a well-integrated DLP system can act in real time: blocking an
email carrying a flagged attachment before it sends, or quarantining a file upload before it
completes. This makes DLP one of the few controls in this course that can stop a loss \emph{during}
the attempt rather than only detecting it afterward. An attacker aware of this will try to blend
exfiltration into normal-looking traffic specifically to avoid triggering it.
\end{example}

\begin{misconception}
\textbf{``A large data transfer is always suspicious.''} Legitimate large transfers -- backups,
sanctioned partner data sharing -- happen constantly. Effective DLP relies on an approved-destination
list and account-specific baselines, not raw transfer size alone.
\end{misconception}

\section{File Integrity Monitoring}

Logs and DLP both watch data in motion; a complementary control watches whether the files
themselves have quietly changed. \textbf{File Integrity Monitoring (FIM)} computes a cryptographic
hash of every file that matters -- binaries, configuration files, web application source -- while
the system is in a known-good state. It stores those hashes as a trusted baseline and periodically
recomputes them to detect drift. Any mismatch means the file changed. If no authorized change ticket
explains it, that mismatch is evidence of compromise -- this is the hash-comparison idea from
Chapter~\ref{ch:5.3} applied directly to forensics rather than to passwords.

\begin{lstlisting}[caption={A hand-rolled SHA-256 FIM baseline (Linux/macOS)}]
# Establish a baseline for the web root
sha256sum -- /var/www/html/*.php /var/www/html/uploads/* > /opt/fim/baseline_webroot.sha256

# Later (e.g. via cron every 15 minutes), verify against the baseline
sha256sum -c /opt/fim/baseline_webroot.sha256 --quiet
# Example output on a compromised host:
# /var/www/html/uploads/profile_pic_884.php: FAILED
# sha256sum: WARNING: 1 computed checksum did NOT match
\end{lstlisting}

\begin{lstlisting}[caption={The equivalent check on Windows with PowerShell}]
# Baseline: compute and store a SHA-256 digest for each watched file
Get-ChildItem C:\inetpub\wwwroot -Recurse -File |
    Get-FileHash -Algorithm SHA256 |
    Export-Csv C:\fim\baseline_webroot.csv -NoTypeInformation

# Later: recompute and compare against the stored baseline
$baseline = Import-Csv C:\fim\baseline_webroot.csv
Get-ChildItem C:\inetpub\wwwroot -Recurse -File | Get-FileHash -Algorithm SHA256 |
    Where-Object { $_.Hash -notin $baseline.Hash } |
    Select-Object Path, Hash
\end{lstlisting}

A hand-rolled hash baseline is instructive but does not scale to enterprise use, which is why
purpose-built FIM tools exist. \textbf{Tripwire}, first released in 1992 by Gene Kim, then a Purdue
graduate student, pioneered the model of a policy file defining watched paths plus a cryptographically
signed database, preventing an attacker who gains root from simply editing the baseline to hide
their tracks.\srccite{5.6}{1} \textbf{AIDE (Advanced Intrusion Detection Environment)} is its
open-source, GPL-licensed replacement.\srccite{5.6}{2} The reference the hashes are compared against
must itself be trustworthy. A baseline stored on the same compromised host an attacker already
controls can simply be rewritten alongside the files it is meant to catch, which is why a signed,
off-host baseline matters as much as the hashing itself.

\subsection{Judging a Detection Method}

Three controls have been introduced so far -- log-based monitoring, DLP, and FIM -- and none of
them is a strictly better or worse option than the others. Judging any detection method means
asking the same five questions: what can it actually see, how much noise does it generate per
genuine finding, and how fast does it fire? What does it cost to operate, and can an attacker who
compromises the host disable it?

Log-based monitoring sees almost anything the application is configured to record, but a busy
server's raw volume means untuned rules drown a real alert in noise. An attacker with code execution
can delete local logs outright, unless they are streamed off-host in real time. DLP sees data only
as it crosses a boundary, and can block a transfer before it completes -- something neither logs nor
FIM can do. Content-based DLP is blind to anything encrypted; volume-based DLP produces false
positives until its baselines are tuned. FIM detects unauthorized file changes cheaply and
precisely, but only after the fact -- it blocks nothing. Its own baseline is as vulnerable to on-host
tampering as any other file, unless stored off-host and signed.

No single method wins on every axis. A mature detection program layers all three specifically
because each control's blind spot is a different control's strength.

\section{Digital Evidence Analysis and Attack Classification}

Detecting that something happened is only half of incident response. Once an alert fires, an
analyst has to answer the harder question of what actually happened and what evidence supports that
conclusion.

\begin{definition}{Attack classification}
The process of collecting multiple indicators of compromise and reasoning about which attack
category they collectively and most strongly support, rather than treating any single indicator as
conclusive on its own.
\end{definition}

An alert is a hypothesis, not a verdict, and this echoes Chapter~\ref{ch:4.4}'s IOC-scoring logic
directly. A spike in failed logins \emph{suggests} brute-force, but confirming it requires checking
whether the failures come from one source targeting many usernames (brute-force) or many sources
targeting one username (credential stuffing from a leaked password list). Those are two genuinely
different attacks that a single 401-count alone cannot distinguish.

\begin{lstlisting}[language=Python, caption={Scoring collected evidence against known attack categories}]
CATEGORY_INDICATORS = {
    "sql_injection": {"sql_metachar_in_request", "database_error_response"},
    "phishing": {"suspicious_attachment", "credential_harvest_url", "urgency_language"},
    "brute_force": {"high_401_volume", "single_source_ip", "many_usernames_targeted"},
}

def classify_attack(observed_indicators):
    """observed_indicators: set of indicator name strings.
       Returns the category with the highest indicator-overlap score,
       or None if no category has any overlap."""
    scores = {cat: len(indicators & observed_indicators)
              for cat, indicators in CATEGORY_INDICATORS.items()}
    best = max(scores, key=scores.get)
    return best if scores[best] > 0 else None
\end{lstlisting}

Evidence supporting a real response decision -- or a legal proceeding -- must be handled with a
documented \textbf{chain of custody}: who collected it, when, how it was stored, and who has touched
it since. An analyst who copies a log to a personal laptop and edits it ``to make it more readable''
has compromised that evidence -- not necessarily because it was tampered with, but because no one
can prove it was not. Proper practice works from read-only copies and hashes the original evidence
at collection time, using exactly the hash-verification idea from Chapter~\ref{ch:5.3}. That way any
later change is immediately detectable.

\begin{figure}[h]
\centering
\begin{tikzpicture}[
  box/.style={draw=kickstartermaroon, thick, rounded corners, minimum width=2.7cm, minimum height=0.9cm, align=center, font=\scriptsize},
]
\node[box, fill=kickstartercream] (collect) at (0,0) {Collect evidence,\\hash immediately};
\node[box] (store) at (3.4,0) {Store read-only\\copy, log custody};
\node[box] (analyze) at (6.8,0) {Analyze a\\working copy};
\node[box, fill=kickstartermaroon!20] (verify) at (10.2,0) {Re-hash original,\\confirm unchanged};
\draw[-{Latex[length=2mm]}] (collect) -- (store);
\draw[-{Latex[length=2mm]}] (store) -- (analyze);
\draw[-{Latex[length=2mm]}] (store) -- (verify);
\node[font=\scriptsize\sffamily, gray, align=center] at (5.1,-1.1)
  {every step after collection works from a copy -- the hashed original is never itself edited};
\end{tikzpicture}
\caption{Chain of custody: the original is hashed at collection and never touched again; every
subsequent step works from a copy, and re-hashing the original at any later point proves nothing
was altered in between.}
\label{fig:5.6-chain-of-custody}
\end{figure}

Industry frameworks like MITRE ATT\&CK give analysts a shared vocabulary, so that ``brute-force''
means the identical specific thing across every team and report.\srccite{5.6}{3} Without it, one
analyst's classification cannot be meaningfully compared or aggregated with another's.

Attack classification is rarely a solo act. A junior analyst collects and documents evidence, a
senior analyst reviews the classification and challenges weak reasoning before it becomes an
official record, and the final classification determines which response playbook a team follows
from there. Honestly communicating uncertainty (``the evidence most strongly supports phishing, but
a malware component cannot be ruled out'') is the mark of a mature analyst, not a weakness in the
analysis.

\begin{misconception}
\textbf{``One suspicious indicator is enough to declare the attack type.''} A single IOC is rarely
conclusive on its own, exactly as in Chapter~\ref{ch:4.4}. Classification should weigh the full set
of collected evidence, not the loudest single signal.
\end{misconception}

\section*{Review Questions}
\begin{reviewquestions}
\item A log shows two failed logins followed by a success for one account within five seconds.
Explain why this pattern alone is not conclusive evidence of an attack, and describe what
additional evidence would change that assessment.
\item Explain why volume/behavior-based DLP detection can catch an exfiltration attempt that
content-based detection completely misses, using an encrypted file transfer as your example.
\item An analyst copies an incident's log file to a personal laptop and edits formatting for
readability before sharing it with a colleague. Explain why this compromises the evidence's chain
of custody even if the analyst made no substantive changes to its content.
\item A SOC receives a single alert: one account exceeded a 401-response threshold. Using the
attack-classification reasoning from this chapter, explain why this alone cannot distinguish
brute-force from credential stuffing, and describe the specific evidence that would resolve the
ambiguity.
\item Explain why logs should stream to a remote, append-only aggregator in real time rather than
being stored only locally on the server that generated them, connecting your answer to what an
attacker with code execution typically tries to do next.
\item An organization stores its file-integrity-monitoring baseline on the same server the baseline
is meant to protect. Explain why an attacker who gains root on that server can defeat the entire
control, and describe what property the baseline's storage location must have to prevent this.
\item Using the chain-of-custody diagram, explain why analysis happens on a working copy rather
than on the hashed original, and what specific guarantee would be lost if an analyst skipped that
step and ran their tools directly against the original evidence file.
\item A junior analyst classifies an incident as phishing, and a senior analyst reviewing the same
evidence says the classification most strongly supports phishing but cannot rule out a malware
component. Explain why the senior analyst's hedged framing reflects better practice than the
junior analyst's confident single-category call, using this chapter's own standard for what makes
a classification credible.
\item Using the five-question framework for judging a detection method, compare FIM to DLP: name
one axis on which FIM is the stronger control, and one axis on which DLP is the stronger control,
citing what this chapter says about each.
\end{reviewquestions}

\begin{answerkey}
\item Two failed logins followed by a success is consistent with an ordinary user mistyping a
password twice, which happens constantly and carries no attack implication by itself. It would
become concerning evidence if the same pattern recurred across many usernames from one source in a
short window (brute-force), or the same username were targeted from many different sources
(credential stuffing). A single instance lacks the correlated volume that gives the pattern
evidentiary weight.
\answerlabel{5.6.A}
\item Content-based detection inspects the data itself for recognizable patterns like a credit-card
number format, but an encrypted file's contents look like random bytes with no such pattern to
match. Content inspection sees nothing suspicious. Volume/behavior-based detection instead watches
the shape of the transfer -- an account that normally moves a few megabytes a day suddenly sending
two gigabytes to an unfamiliar destination is anomalous regardless of what the encrypted payload
actually contains. It flags the transfer purely on its behavioral signature.
\answerlabel{5.6.B}
\item Chain of custody requires being able to prove evidence was not altered after collection, not
merely that it was not altered in substance. Once a copy exists on an uncontrolled personal device
and has been edited at all, no one can prove from the evidence itself that only formatting changed
and not content. The provable guarantee is gone even if the actual change was harmless, which is why
proper practice works from read-only copies and hashes the original at collection time instead.
\answerlabel{5.6.E}
\item A high count of 401 responses from a single account could mean an attacker is guessing that
one account's password from many sources (credential stuffing). It could also mean the alert is
actually describing brute-force targeting many accounts and this is just one of them. The raw count
alone does not reveal whether the failures are concentrated on one source IP against many usernames,
or distributed across many source IPs against this one username. Correlating the failed attempts by
source IP and target username resolves the ambiguity: one IP hitting many usernames is brute-force,
many IPs hitting one username is credential stuffing.
\answerlabel{5.6.E}
\item An attacker who achieves code execution commonly tries to erase local logs to cover their
tracks. If those logs exist only on the compromised server, that deletion succeeds completely and
destroys the forensic record along with the evidence of the attack itself. Streaming logs to a
separate, append-only, remote aggregator in real time means the attacker's local deletion cannot
reach the copies that already left the machine, preserving the record regardless of what happens to
the original.
\answerlabel{5.6.B}
\item File Integrity Monitoring only detects tampering by comparing current file hashes against a
trusted baseline. If the baseline sits on the same host, an attacker with root can recompute and
overwrite the stored baseline hashes to match their tampered files, and every subsequent check
reports a clean match because the reference itself was falsified. The baseline's storage must be
off-host or otherwise outside the attacker's reach once they have compromised the monitored server
-- for example, a separate, access-controlled system, or a cryptographically signed database like
Tripwire's that requires a separate credential to update. Compromising the monitored host alone is
not sufficient to also rewrite the reference it is checked against.
\answerlabel{5.6.D}
\item Analysis happens on a working copy so that whatever tools, searches, or annotations the
analysis requires can never alter the original evidence. If something needs to be tried again
differently, or a mistake is made while analyzing, the untouched original is still there to copy
from again. If an analyst instead ran tools directly against the original, any tool that modifies
timestamps, reformats content, or otherwise touches the file during analysis would break the
provable guarantee that the evidence is unchanged since collection -- the entire point of having
hashed it in the first place. The loss is not necessarily that the content changes, but that no one
could prove it did not.
\answerlabel{5.6.E}
\item This chapter's standard is that classification should weigh the full set of collected
evidence rather than the loudest single signal, and that honestly communicating uncertainty is the
mark of a mature analyst, not a weakness. A confident single-category call implies a level of
certainty the underlying evidence may not actually support. Explicitly naming what remains
unresolved instead gives whoever acts on the classification an accurate picture of the actual
evidentiary strength -- including the possibility that the response playbook may need to account for
a scenario the primary classification does not cover.
\answerlabel{5.6.E}
\item FIM is the stronger control on cost and noise: once a baseline exists, comparing hashes is
cheap to run and generates almost no false positives, since a mismatch is either an authorized
change or a real compromise, unlike volume-based DLP, which needs tuned baselines and an
approved-destination list to avoid raising them. DLP is the stronger control on speed and
prevention: a well-integrated DLP system can block a transfer while it is still happening, while FIM
only detects a change that has already occurred and never prevents the write in the first place.
\answerlabel{5.6.C}
\end{answerkey}

\section*{Sources}
\begin{enumerate}[label=\lbrack\arabic*\rbrack, leftmargin=2.2em]
\item \srcitem{5.6}{1}Gene H.\ Kim and Eugene H.\ Spafford, ``The Design and Implementation of
Tripwire: A File System Integrity Checker,'' Purdue University Computer Science Technical Report
CSD-TR-93-071 (1993); Tripwire, Inc., ``The Past, Present, and Future of File Integrity Monitoring,''
\emph{The State of Security}.
\item \srcitem{5.6}{2}``Advanced Intrusion Detection Environment (AIDE),'' project homepage,
aide.github.io.
\item \srcitem{5.6}{3}MITRE Corporation, ``MITRE ATT\&CK\textregistered,'' attack.mitre.org.
\end{enumerate}
